\documentclass[aspectratio=169]{beamer}
\usepackage[norsk]{babel}
\usepackage{booktabs}
\usepackage[T1]{fontenc}
\usepackage[utf8]{inputenc}
\usepackage{graphicx,epstopdf,tikz, ifthen}
\usepackage{siunitx, enumerate, xcolor, hyperref}
\usepackage{ enumerate, xcolor}
\usepackage[export]{adjustbox}
\usepackage{circuitikz}
%\usepackage{enumitem}
%\usepackage{enumitem, siunitx}
\usepackage{multicol}
%\usepackage{mathtools}
\usetikzlibrary{decorations.pathreplacing, calc, decorations.pathmorphing, patterns}

\usepackage{minted}

\usetheme[slogan=norsk,style=vertical, mathfont=serif]{NTNU}
\setbeamertemplate{section in toc}[circle]

\title[rl]{Mappevurdering \\ IIRA2001}
\date{14. November 2024}
\author{Jonas J. Harang}

%SIunitx
\let\DeclareUSUnit\DeclareSIUnit
\let\US\SI
\let\us\si

\DeclareUSUnit\gallon{gal}
\DeclareUSUnit\inch{in}
\DeclareUSUnit\psi{psi}

\DeclareSIUnit\atm{atm}
\DeclareSIUnit\bar{bar}
\DeclareSIUnit\mmhg{mm Hg}

%%--- tikzcolor for ntnu
\definecolor{NTNUblue}{HTML}{00509e}
\definecolor{NTNUlightblue}{HTML}{6096d0}
\definecolor{NTNUorange}{HTML}{ef8114}
\definecolor{NTNUbrown}{HTML}{cfb887}
\definecolor{NTNUred}{HTML}{b01b81}
\definecolor{NTNUgrey}{HTML}{bebebe}
\definecolor{NTNUcyan}{HTML}{3cbfbe}
\definecolor{NTNUviolet}{HTML}{482776}



%-----------Beamer jonas
%\setbeamertemplate{subsection in toc}[subsections numbered]
\setbeamertemplate{subsection in toc}[subsections unnumbered]
\makeatletter
\newcommand{\SetSectionNumber}[1]{%
\beamer@tocsectionnumber=\numexpr#1-1}
\makeatother

\begin{document}

\begin{frame}
\maketitle
\end{frame}

\begin{frame}{Info om vurdering}
\end{frame}

\begin{frame}{Viktige avklaringer: Individuell vurdering}
   \huge \color{red}{Mappevurderingen og innleveringen er Individuell}
   \small 

   \color{black}
   \begin{columns}
      \begin{column}{0.6\textwidth}
         \begin{block}{Det betyr}
            \begin{enumerate}[$\bullet$]
               \item {Du leverer et individuelt arbeid -- din egen besvarelse som \textit{du} har tilvirket}
               \item {Dere går IKKE sammen i grupper og leverer samme innlevering (katastrofe)}
               \item {At man samarbeider med hverandre med kodingen slik at koden ser lik ut går stort sett fint}
               \item {At man samarbeider med hverandre og fordeler arbeidet i en gruppe slik at koden er identisk går \textit{ikke} fint}
            \end{enumerate}
         \end{block}
      \end{column}
      
      \begin{column}{0.4\textwidth}
         \begin{block}{Hvor strenge er man?}
            \begin{enumerate}[$\bullet$]
               \item {Det er greit at man deler litt kode med hverandre}
               \item {Det er greit at man klipper og limer litt fra forelesning}
               \item {Vi bygger egen kode basert på eksempler man har lært}
               \item {I første del av mappen er det naturlig at programmene blir veldig like, det er helt greit}
            \end{enumerate}
         \end{block}
      \end{column}
   \end{columns}
\end{frame}

\begin{frame}{Viktige avklaringer: Individuell vurdering}
   \begin{columns}
      
      \begin{column}{0.5\textwidth}
         \begin{block}{Mappe del 1:}
            Dersom dere har samarbeidet tett i del 1:
            \begin{itemize}
               \item {Pass på at rapport er individuell}
               \item {Bruk litt andre tall/input slik at eventuell diskusjon blir annerledes}
               \item {Ta kontakt med faglærer om dere merker at det er vanskelig å skille løsningene deres}
               \item {Vi forventer som sagt at løsningne er like til del 1}
            \end{itemize}
         \end{block}
      \end{column}
      
      \begin{column}{0.5\textwidth}
         \begin{enumerate}[$\bullet$]
            \item {I del 2 av mappen er dere frie til å velge datakilde og datasett}
            \item {Dere velger også egen problemstilling eller vinkling på besvarelsen}
            \item {Da blir det viktigere å passe på at dere leverer en individuell besvarelse}
            \item {Det er sunt og bra for læringsprosessen å samarbeide med hverandre, slik at vi vil at denne muligheten skal være åpen}
         \end{enumerate}
      \end{column}
   \end{columns}
\end{frame}

\begin{frame}{Viktige avklaringer: Individualisert besvarelse del 2}
   Dersom man jobber tett til besvarelsene av del 2:
   \small
   \begin{columns}
      \begin{column}{0.5\textwidth}
         \begin{block}{Dere har kanskje:}
            \begin{enumerate}[$\bullet$]
               \item {Jobbet med samme dataset}
               \item {Har samme problemstilling eller vinkling}
            \end{enumerate}
         \end{block}
         \begin{block}{Samme dataset}
           \begin{enumerate}[$\bullet$]
               \item {Det er greit at dere løser det å lese inn, vaske og prepare data likt}
               \item {Dere burde da ha ulike problemstillinger eller ulike formål med dataene}
               \item {Tommelfingerregel: Blir resultat og diskusjon annerledes er det greit}
            \end{enumerate}
         \end{block}
      \end{column}
      \begin{column}{0.5\textwidth}
         
         \begin{block}{Problemstilling/Vinkling}
            \begin{enumerate}[$\bullet$]
               \item {Kanskje er dere flere som ønsker å undersøke korrelasjon, utforske og visualisere osv.}
               \item {Det er greit at dere løser det å lese inn/preparere data, lager plot, kjører analyse osv. likt}
               \item {Dere burde da bruke forskjellige datagrunnlag eller datasett}
               \item {Tommelfingerregel: Blir resultat og diskusjon annerledes er det greit}
            \end{enumerate}
         \end{block}
      \end{column}
   \end{columns}
\end{frame}

\begin{frame}{Hjeeeelp, vi har misforstått}
   \huge
   Ta kontakt med faglærerer dersom dere er bekymret for om dere har overskridd grensen til plagiat asap, så planlegger vi en løsning sammen
\end{frame}


\begin{frame}{Viktige avklaringer: Innlevering}
   \begin{itemize}
      \item {Mappen leveres på inspera}
      \item {Både del 1 og del 2 leveres inn}
      \item {Levering under workshop og medstudentvurdering på iirmoodle telles ikke som en del av mappebesvarelsen}
      \item {Det åpnes for innlevering i inspera 12.09 klokken 09:00, og det stenges 13.12 klokken 12:00}
      \item {Info om hvordan man får levert legges ut på blackboard}
   \end{itemize}
\end{frame}

\begin{frame}{Viktige avklaringer: Oppgavebeskrivelse del 2}
   \begin{itemize}
      \item {Oppgaveskrivelsen til del 2 av mappen sier den er \textit{tentativ}}
      \item {Det gjaldt i all hovedsak tips og eksemplene}
      \item {\textit{Selve} oppgaven blir stående slik den er}
   \end{itemize}
\end{frame}

\begin{frame}{Viktige avklaringer: Bruk av KI}
   \begin{itemize}
      \item {Flittig bruk av KI-til brainstorming og som studass er gode læringsstrategier}
      \item {Til mappen leveres KI-deklarasjonsskjema}
      \item {Eventuell koden i innleveringen som er generert av feks ChatGPT merkes med kommentarer, så går dette fint}
   \end{itemize}
\end{frame}

\begin{frame}[fragile]{Eksempel, KI-merking}
   \tiny
   \begin{minted}{python}
import numpy as np

def incmatrix(genl1,genl2):
    m = len(genl1)
    n = len(genl2)
    M = None #to become the incidence matrix
    VT = np.zeros((n*m,1), int)  #dummy variable

    #compute the bitwise xor matrix
    M1 = bitxormatrix(genl1)
    M2 = np.triu(bitxormatrix(genl2),1)

    #Følgende blokk er generert av KI ----------->
    for i in range(m-1):
        for j in range(i+1, m):
            [r,c] = np.where(M2 == M1[i,j])
            for k in range(len(r)):
                VT[(i)*n + r[k]] = 1;
                VT[(i)*n + c[k]] = 1;
                VT[(j)*n + r[k]] = 1;
                VT[(j)*n + c[k]] = 1;

                if M is None:
                    M = np.copy(VT)
                else:
                    M = np.concatenate((M, VT), 1)

                VT = np.zeros((n*m,1), int)
    #<------------- Generert av KI

    return M
   \end{minted}
\end{frame}

\begin{frame}{Sensur}
   \huge
   Oppdatert sensurveiledning er tilgjengelig på blackboard
   
   Uenigheter eller forslag til endringer kan tas opp gjennom referansegruppen
\end{frame}

\begin{frame}{Endringer}
   \begin{itemize}
      \item {Grunnet tid og helseproblematikken i semester ser jeg ikke mulighet til å gi faglærers tilbakemelding til del 1 av mappen}
      \item {Vi skal gå gjennom 2 besvarelser fra i fjor straks}
      \item {Er det ønskelig å kjøre en ny workshop med medstudentvurdering til mappe 2? (Uten gruppeinnlevering)}
   \end{itemize}
\end{frame}

\begin{frame}{Løsningseksempel}
   \begin{itemize}
      \item {Har fått tilbakemelding om vansker med å henge med, og bekymringer rundt vanskelighetsgrad}
      \item {Etter det jeg har sendt i klasserommet er jeg overhode ikke bekymret (moderat bekymret for de ca. 50\% jeg ikke ser)}
      \item {Vi går gjennom Løsningseksempel fra i fjor for å klargjøre forventninger til sensur}
   \end{itemize}
\end{frame}


\end{document}

